% GENERATED FILE. Edit canonical lessons, then run npm run generate:books.
% canonical-source-hash: fnv1a64:ce0c96eef5bb3462
% canonical-lessons: PT-C488-repouso, PT-C488-curar, PT-C488-cura, PT-C488-recuperar, PT-C488-recuperacao

\chapter{Ficar bom}
\label{ch:pt-ficar-bom}
\hlchaptermodality{488}

\begin{chapteropening}
\textbf{By the end of this chapter:} \emph{I can talk about getting better after an illness.}

Complete the last lesson of chapter 488: I can talk about getting better after an illness.
\end{chapteropening}

\section[o repouso]{o repouso — rest}

\label{lesson:PT-C488-repouso}



\begin{quote}
Before the new one: say the Portuguese for to prescribe, then the Portuguese for treatment.
\end{quote}

\subsection*{You'll want to know: o repouso}
\textbf{o repouso} — \textquotedblleft{}rest\textquotedblright{}.

\begin{quote}
O médico recomendou uma semana de repouso.
\end{quote}

\begin{quote}
\emph{The doctor recommended a week of rest.}
\end{quote}

\emph{Estar em repouso} — to be resting, on doctor's orders.

\begin{cousinweb}
\emph{repouso} comes from \emph{repousar} \textquotedblleft{}to rest\textquotedblright{} (EtymDB), inherited from Late Latin \emph{repausāre} \textquotedblleft{}to lay at rest\textquotedblright{}, built on \emph{pausāre} \textquotedblleft{}to pause\textquotedblright{}.
\end{cousinweb}

\begin{grammarlens}[title={What you've built}]
One more word for telling, explaining and describing at B1.
\end{grammarlens}

\subsection*{Your turn}
\begin{itemize}
\raggedright
  \item \emph{Say it:} \emph{o repouso}
  \item \emph{Say it:} \emph{o repouso}, once more
  \item \emph{Say it:} a sentence of your own with \emph{o repouso}
  \item \emph{Recall:} say the Portuguese for to prescribe, then the Portuguese for treatment, then say \emph{o repouso} again
\end{itemize}

\begin{tcolorbox}[breakable,colback=teal!4,colframe=teal!35!black,title={Before you move on}]
What does o repouso mean? (\textquotedblleft{}Rest\textquotedblright{}.) Say it once more.
\end{tcolorbox}
\section[curar]{curar — to cure, to heal}

\label{lesson:PT-C488-curar}



\begin{quote}
Before the new one: say the Portuguese for treatment, then the Portuguese for rest.
\end{quote}

\subsection*{You'll want to know: curar}
\textbf{curar} — \textquotedblleft{}to cure, to heal\textquotedblright{}.

\begin{quote}
O xarope não curou a tosse, mas ajudou.
\end{quote}

\begin{quote}
\emph{The syrup did not cure the cough, but it helped.}
\end{quote}

\emph{Curar-se} — to get well.

\begin{cousinweb}
\emph{curar} is inherited from Latin \emph{cūrāre} \textquotedblleft{}to take care of\textquotedblright{} (EtymDB), built on \emph{cūra} \textquotedblleft{}care\textquotedblright{}.
\end{cousinweb}

\begin{grammarlens}[title={What you've built}]
One more word for telling, explaining and describing at B1.
\end{grammarlens}

\subsection*{Your turn}
\begin{itemize}
\raggedright
  \item \emph{Say it:} \emph{curar}
  \item \emph{Say it:} \emph{curar}, once more
  \item \emph{Say it:} a sentence of your own with \emph{curar}
  \item \emph{Recall:} say the Portuguese for treatment, then the Portuguese for rest, then say \emph{curar} again
\end{itemize}

\begin{tcolorbox}[breakable,colback=teal!4,colframe=teal!35!black,title={Before you move on}]
What does curar mean? (\textquotedblleft{}To cure, to heal\textquotedblright{}.) Say it once more.
\end{tcolorbox}
\section[a cura]{a cura — cure}

\label{lesson:PT-C488-cura}



\begin{quote}
Before the new one: say the Portuguese for rest, then the Portuguese for to cure.
\end{quote}

\subsection*{You'll want to know: a cura}
\textbf{a cura} — \textquotedblleft{}cure\textquotedblright{}.

\begin{quote}
Ainda não há cura para esta doença.
\end{quote}

\begin{quote}
\emph{There is still no cure for this disease.}
\end{quote}

Also a stay for health: \emph{uma cura de águas}.

\begin{cousinweb}
\emph{cura} is inherited from Latin \emph{cūra} \textquotedblleft{}care, concern\textquotedblright{} (EtymDB).
\end{cousinweb}

\begin{grammarlens}[title={What you've built}]
One more word for telling, explaining and describing at B1.
\end{grammarlens}

\subsection*{Your turn}
\begin{itemize}
\raggedright
  \item \emph{Say it:} \emph{a cura}
  \item \emph{Say it:} \emph{a cura}, once more
  \item \emph{Say it:} a sentence of your own with \emph{a cura}
  \item \emph{Recall:} say the Portuguese for rest, then the Portuguese for to cure, then say \emph{a cura} again
\end{itemize}

\begin{tcolorbox}[breakable,colback=teal!4,colframe=teal!35!black,title={Before you move on}]
What does a cura mean? (\textquotedblleft{}Cure\textquotedblright{}.) Say it once more.
\end{tcolorbox}
\section[recuperar]{recuperar — to recover}

\label{lesson:PT-C488-recuperar}



\begin{quote}
Before the new one: say the Portuguese for to cure, then the Portuguese for cure.
\end{quote}

\subsection*{You'll want to know: recuperar}
\textbf{recuperar} — \textquotedblleft{}to recover\textquotedblright{}.

\begin{quote}
Recuperei depressa depois da operação.
\end{quote}

\begin{quote}
\emph{I recovered quickly after the operation.}
\end{quote}

\emph{Recuperar de} — to recover from; \emph{recuperar o tempo perdido} — to make up for lost time.

\begin{cousinweb}
\emph{recuperar} is a learned borrowing from Latin \emph{recuperāre} \textquotedblleft{}to get back, to regain\textquotedblright{} (EtymDB).
\end{cousinweb}

\begin{grammarlens}[title={What you've built}]
One more word for telling, explaining and describing at B1.
\end{grammarlens}

\subsection*{Your turn}
\begin{itemize}
\raggedright
  \item \emph{Say it:} \emph{recuperar}
  \item \emph{Say it:} \emph{recuperar}, once more
  \item \emph{Say it:} a sentence of your own with \emph{recuperar}
  \item \emph{Recall:} say the Portuguese for to cure, then the Portuguese for cure, then say \emph{recuperar} again
\end{itemize}

\begin{tcolorbox}[breakable,colback=teal!4,colframe=teal!35!black,title={Before you move on}]
What does recuperar mean? (\textquotedblleft{}To recover\textquotedblright{}.) Say it once more.
\end{tcolorbox}
\section[a recuperação]{a recuperação — recovery}

\label{lesson:PT-C488-recuperacao}



\begin{quote}
Before the new one: say the Portuguese for cure, then the Portuguese for to recover.
\end{quote}

\subsection*{You'll want to know: a recuperação}
\textbf{a recuperação} — \textquotedblleft{}recovery\textquotedblright{}.

\begin{quote}
A recuperação da perna partida demorou três meses.
\end{quote}

\begin{quote}
\emph{The recovery from the broken leg took three months.}
\end{quote}

Built on \emph{recuperar}.

\begin{cousinweb}
\emph{recuperação} is built on \emph{recuperar}, a learned borrowing from Latin \emph{recuperāre} \textquotedblleft{}to regain\textquotedblright{} (EtymDB); compare Latin \emph{recuperātiō}.
\end{cousinweb}

\begin{grammarlens}[title={What you've built}]
One more word for telling, explaining and describing at B1.
\end{grammarlens}

\subsection*{Your turn}
\begin{itemize}
\raggedright
  \item \emph{Say it:} \emph{a recuperação}
  \item \emph{Say it:} \emph{a recuperação}, once more
  \item \emph{Say it:} a sentence of your own with \emph{a recuperação}
  \item \emph{Recall:} say the Portuguese for cure, then the Portuguese for to recover, then say \emph{a recuperação} again
\end{itemize}

\begin{tcolorbox}[breakable,colback=teal!4,colframe=teal!35!black,title={Before you move on}]
What does a recuperação mean? (\textquotedblleft{}Recovery\textquotedblright{}.) Say it once more.
\end{tcolorbox}
